\documentclass[../../main.tex]{subfiles}
\begin{document}

\begin{center}
  {\fontsize{48pt}{58pt}\selectfont \textbf{1992}\normalsize\textbf{年度}}
\end{center}

\begin{table}[htb]
  \centering
  \caption{出題テーマと難易度}
  \begin{tabular}{|c|p{9cm}|c|}
    \hline
    \textbf{問題番号} & \textbf{テーマ} & \textbf{難易度} \\
    \hline
    第1問 & 正八面体の平面切断と頂点座標 & \\
    \hline
    第2問 & 三角方程式と整数係数の決定 & \\
    \hline
    第3問 & 外心から下ろした垂線と三角形の角 & \\
    \hline
    第4問 & サイコロの目の積の倍数確率 & \\
    \hline
    第5問 & 回転放物面への接平面で囲む立体 & \\
    \hline
    第6問 & ピタゴラス数の漸化式と比 & \\
    \hline
  \end{tabular}
\end{table}

\begin{table}[htb]
  \centering
  \caption{解答}
  \renewcommand{\arraystretch}{1.4}
  \begin{tabular}{|c|c|p{9cm}|}
    \hline
    \textbf{問題番号} & \textbf{小問} & \textbf{答え} \\
    \hline
    第1問 & -- & 切り口は六角形．頂点は$\left(\pm\dfrac{a}{a+b},\mp\dfrac{b}{a+b},0\right)$，$\left(\pm\dfrac{a}{a+c},0,\mp\dfrac{c}{a+c}\right)$，$\left(0,\pm\dfrac{b}{b+c},\mp\dfrac{c}{b+c}\right)$ \\
    \hline
    第2問 & (1) & $\theta=\dfrac{\pi}{5}$ \\
    \cline{2-3}
    & (2) & $(m,n,\theta)=(0,0,\pi/3),(0,1,\pi/4),(0,2,\pi/6),(1,0,\pi/5)$ \\
    \hline
    第3問 & (1) & $5\overrightarrow{OA}+4\overrightarrow{OB}+3\overrightarrow{OC}=\vec{0}$ \\
    \cline{2-3}
    & (2) & $\angle A=\dfrac{\pi}{4}$ \\
    \hline
    第4問 & (1) & $p_n=1-\left(\dfrac{2}{3}\right)^n$ \\
    \cline{2-3}
    & (2) & $q_n=1+\left(\dfrac{1}{3}\right)^n-\left(\dfrac{1}{2}\right)^n-\left(\dfrac{2}{3}\right)^n$ \\
    \hline
    第5問 & (1) & $V(t)=\dfrac{(t^2+1)^3}{3t^2}$ \\
    \cline{2-3}
    & (2) & $t=\dfrac{1}{\sqrt{2}}$で最小値$\dfrac{9}{4}$ \\
    \hline
    第6問 & (1) & $a_n^2+b_n^2=c_n^2$（証明問題） \\
    \cline{2-3}
    & (2) & $c_n>0$かつ$c_n\ge c_{n+1}$（証明問題） \\
    \cline{2-3}
    & (3) & $a_m:b_m:c_m=3:4:5$または$4:3:5$ \\
    \hline
  \end{tabular}
\end{table}

\paragraph{試験形態}
6問．試験時間は要確認．

\end{document}
